\documentclass[11pt]{article}

\usepackage[a4paper,margin=1in]{geometry}
\usepackage[T1]{fontenc}
\usepackage[utf8]{inputenc}
\usepackage{lmodern}
\usepackage{microtype}
\usepackage{amsmath}
\usepackage{booktabs}
\usepackage{enumitem}
\usepackage{hyperref}
\usepackage{xcolor}
\usepackage{siunitx}

\hypersetup{
  colorlinks=true,
  linkcolor=blue!50!black,
  urlcolor=blue!50!black,
  citecolor=blue!50!black
}

\title{LZ $0\nu\beta\beta$ Sensitivity: Key Parameters and Background Summary}
\author{Extracted from arXiv:1912.04248v2}
\date{\today}

\begin{document}
\maketitle

\begin{abstract}
This note summarizes the key detector assumptions, background model,
and projected sensitivity of the LUX-ZEPLIN (LZ) experiment to
neutrinoless double beta decay of $^{136}$Xe, as reported in
arXiv:1912.04248v2. LZ projects a median 90\% CL exclusion
sensitivity of $1.06 \times 10^{26}$~years for the $^{136}$Xe
$0\nu\beta\beta$ half-life after 1000~live-days, with a total
expected background of $\sim$35.6 events in the $\pm 1\sigma$ ROI
and inner 967~kg fiducial volume.
\end{abstract}

% ==================================================================
\section{Detector Parameters}
% ==================================================================

LZ is a two-phase liquid xenon TPC with 7~tonnes of active LXe,
yielding 623~kg of $^{136}$Xe at 8.9\% natural abundance.  The TPC
drift region spans 145.6~cm (cathode to gate) with an inner diameter
of 145.6~cm.  Two PMT arrays (253 top, 241 bottom) detect
scintillation (S1) and electroluminescence (S2) signals.

The detector is surrounded by an LXe skin veto ($\sim$2~tonnes,
100~keV threshold over $>$95\% of the volume), an outer detector
(Gd-doped liquid scintillator, 100~keV threshold, 1~$\mu$s
coincidence window), and a water shield.

% ==================================================================
\section{Key Assumptions}
% ==================================================================

Table~\ref{tab:assumptions} collects the analysis assumptions that
drive the sensitivity.

\begin{table}[h]
\centering
\caption{Key analysis assumptions for the LZ $0\nu\beta\beta$ search.}
\label{tab:assumptions}
\begin{tabular}{@{}ll@{}}
\toprule
Parameter & Value \\
\midrule
$Q_{\beta\beta}$ ($^{136}$Xe) & $2457.83 \pm 0.37$~keV \\
Energy resolution ($\sigma/E$ at $Q_{\beta\beta}$) & 1.0\% (conservative) \\
Predicted resolution (NEST + LZ baseline) & 0.88\% \\
XENON1T demonstrated resolution & 0.79\% \\
ROI ($\pm 1\sigma$ around $Q_{\beta\beta}$) & $2433.3 < E < 2482.4$~keV \\
Extended energy range (PLR fit) & $2000 < E < 2700$~keV \\
Minimum vertex separation ($\Delta z$) & 3~mm \\
Signal efficiency (after SS + veto cuts) & 80\% \\
\midrule
Inner fiducial volume (cut-and-count) & 967~kg ($26 < z < 96$~cm, $r < 39$~cm) \\
Extended fiducial volume (PLR) & 5600~kg ($2 < z < 132.6$~cm, $r < 68.8$~cm) \\
Live time & 1000~days \\
$^{136}$Xe exposure (natural abundance) & 1360~kg$\cdot$years \\
\midrule
Skin veto threshold & 100~keV \\
Outer detector veto threshold & 100~keV \\
Veto coincidence window & 1~$\mu$s \\
\bottomrule
\end{tabular}
\end{table}

% ==================================================================
\section{Background Model}
% ==================================================================

\subsection{Dominant gamma lines near $Q_{\beta\beta}$}

Two gamma lines dominate the background near $Q_{\beta\beta}$:

\begin{itemize}
\item $^{208}$Tl at 2614.5~keV ($^{232}$Th chain, BR = 35.9\%):
  $\sim$157~keV above $Q_{\beta\beta}$.  Partially rejected by
  single-scatter cuts (the coincident $\geq$583~keV companion gamma
  tags multi-site events) and by skin/OD veto.

\item $^{214}$Bi at 2447.7~keV ($^{238}$U chain, BR = 1.5\%):
  $\sim$10~keV below $Q_{\beta\beta}$.  Cannot be separated from the
  signal by energy resolution alone.  The low branching ratio is
  fortunate.
\end{itemize}

\subsection{Backgrounds in the $\pm 1\sigma$ ROI (1000~days, inner 967~kg)}

Table~\ref{tab:backgrounds} reproduces the complete background budget
from Table~I of the LZ $0\nu\beta\beta$ paper.

\begin{table}[h]
\centering
\caption{Expected background counts in the $\pm 1\sigma$ ROI
  ($2433$--$2482$~keV) and inner 967~kg fiducial volume for
  1000~live-days, assuming 1.0\% energy resolution at $Q_{\beta\beta}$
  and 3~mm single-scatter rejection in $z$.  Activities marked
  $^\dagger$ are upper limits.}
\label{tab:backgrounds}
\sisetup{round-mode=places, round-precision=2}
\begin{tabular}{@{}lS[table-format=5.0]S[table-format=4.1]S[table-format=4.1]S[table-format=2.2]S[table-format=2.2]S[table-format=2.2]@{}}
\toprule
{Component} & {Mass} & {$^{238}$U-late} & {$^{232}$Th-late} &
  {Counts} & {Counts} & {Total} \\
 & {(kg)} & {(mBq/kg)} & {(mBq/kg)} &
  {from $^{238}$U} & {from $^{232}$Th} & {Counts} \\
\midrule
TPC PMTs          &   91.9 &   3.22 &   1.61 & 2.95 & 0.10 & 3.05 \\
TPC PMT bases     &    2.80 &  75.9  &  33.1  & 1.52 & 0.03 & 1.55 \\
TPC PMT structures&  166   &   1.60 &   1.06 & 2.65 & 0.12 & 2.77 \\
TPC PMT cables    &   88.7 &   4.31 &   0.82 & 1.44 & 0.19 & 1.63 \\
Skin PMTs + bases &    8.59 &  46.0  &  14.9  & 0.75 & 0.02 & 0.78 \\
PTFE walls        &  184   &   0.04 &   0.01 & 0.39 & 0.00 & 0.39 \\
TPC sensors       &    5.02 &   5.82 &   1.88 & 1.19 & 0.00 & 1.19 \\
Field grids + holders & 89.1 & 2.63 &   1.46 & 0.62 & 0.11 & 0.73 \\
FC resistors      &    0.06 & 1350   & 2010   & 2.63 & 0.03 & 2.65 \\
FC rings$^\dagger$ &  93.0 &   0.35 &   0.24 & 0.82 & 0.00 & 0.82 \\
Ti cryostat$^\dagger$ & 2590 & 0.08 &  0.22  & 1.30 & 0.20 & 1.49 \\
Cryo.\ insulation$^\dagger$ & 13.8 & 11.1 & 7.79 & 0.90 & 0.04 & 0.94 \\
Outer detector$^\dagger$ & 22900 & 4.71 & 3.73 & 1.70 & 1.08 & 2.79 \\
Other components  &  438   &   1.83 &   1.65 & 2.10 & 0.31 & 2.41 \\
\midrule
\textbf{Det.\ components subtotal} & & & & \textbf{21.0} & \textbf{2.32} & \textbf{23.3} \\
\midrule
Cavern walls      & 29000  &   3.21 &   8.41 &  {---} & {---} & 11.6 \\
$^{137}$Xe (neutron-induced) & & & & & & 0.28 \\
Internal $^{222}$Rn & & & & & & 0.45 \\
$^{136}$Xe $2\nu\beta\beta$ & & & & & & 0.01 \\
$^{8}$B solar $\nu$ & & & & & & 0.03 \\
\midrule
\textbf{Total}    &        &        &        & \textbf{24.2} & \textbf{10.7} & \textbf{35.6} \\
\bottomrule
\end{tabular}
\end{table}

\subsection{Detailed background rates per year}

Table~\ref{tab:rates} converts the 1000-day counts from
Table~\ref{tab:backgrounds} to annual rates per component and per
isotope chain (scaling factor $365.25/1000 = 0.36525$).

\begin{table}[h]
\centering
\caption{Annual background rates (counts/year) in the $\pm 1\sigma$
  ROI and inner 967~kg volume, broken down by component and decay
  chain.  Derived from Table~\ref{tab:backgrounds}.}
\label{tab:rates}
\begin{tabular}{@{}lS[table-format=1.3]S[table-format=1.3]S[table-format=1.3]@{}}
\toprule
{Component} & {$^{238}$U-late} & {$^{232}$Th-late} & {Total} \\
 & {(cts/yr)} & {(cts/yr)} & {(cts/yr)} \\
\midrule
TPC PMTs             & 1.077 & 0.037 & 1.114 \\
TPC PMT bases        & 0.555 & 0.011 & 0.566 \\
TPC PMT structures   & 0.968 & 0.044 & 1.012 \\
TPC PMT cables       & 0.526 & 0.069 & 0.596 \\
Skin PMTs + bases    & 0.274 & 0.007 & 0.285 \\
PTFE walls           & 0.142 & 0.000 & 0.142 \\
TPC sensors          & 0.435 & 0.000 & 0.435 \\
Field grids + holders& 0.226 & 0.040 & 0.267 \\
FC resistors         & 0.961 & 0.011 & 0.968 \\
FC rings             & 0.300 & 0.000 & 0.300 \\
Ti cryostat          & 0.475 & 0.073 & 0.544 \\
Cryo.\ insulation    & 0.329 & 0.015 & 0.343 \\
Outer detector       & 0.621 & 0.395 & 1.019 \\
Other components     & 0.767 & 0.113 & 0.881 \\
\midrule
\textbf{Det.\ components subtotal} & \textbf{7.67} & \textbf{0.85} & \textbf{8.51} \\
\midrule
Cavern walls ($^{238}$U + $^{232}$Th) & 1.17 & 3.06 & 4.24 \\
\midrule
\textbf{External subtotal} & \textbf{8.84} & \textbf{3.91} & \textbf{12.75} \\
\midrule
$^{137}$Xe (neutron-induced) & {---} & {---} & 0.102 \\
Internal $^{222}$Rn          & {---} & {---} & 0.164 \\
$^{136}$Xe $2\nu\beta\beta$  & {---} & {---} & 0.004 \\
$^{8}$B solar $\nu$          & {---} & {---} & 0.011 \\
\midrule
\textbf{Internal subtotal}   & {---} & {---} & \textbf{0.28} \\
\midrule[\heavyrulewidth]
\textbf{Grand total}         & \textbf{8.84} & \textbf{3.91} & \textbf{13.01} \\
\bottomrule
\end{tabular}
\end{table}

Key observations from the annual rates:
\begin{itemize}
\item $^{238}$U-late dominates over $^{232}$Th-late by a factor of
  $\sim$9 for detector components, because the $^{214}$Bi 2448~keV
  line falls directly under $Q_{\beta\beta}$ while the $^{208}$Tl
  2615~keV line is partially rejected by veto and single-scatter cuts.

\item For cavern walls the ratio reverses: $^{232}$Th contributes
  $\sim$2.6$\times$ more than $^{238}$U, reflecting both the higher
  $^{232}$Th activity in the rock (8.41 vs.\ 3.21~Bq/kg) and the
  penetrating power of the 2615~keV line through the shielding.

\item The top three detector-component contributors per year are:
  TPC PMTs (1.11), outer detector (1.02), and PMT structures (1.01),
  followed closely by FC resistors (0.97).

\item Internal backgrounds ($^{222}$Rn + $^{137}$Xe +
  $2\nu\beta\beta$ + $^{8}$B) contribute only 0.28~counts/year
  ($\sim$2\% of total), making the search overwhelmingly
  externally-background-limited.
\end{itemize}

\subsection{Background composition}

The dominant contributions to the total 35.6 counts in 1000~days are:

\begin{enumerate}
\item \textbf{Detector components} (23.3 counts, 65\%): dominated by
  TPC PMTs (3.05), PMT structures (2.77), FC resistors (2.65), other
  components (2.41), and PMT cables (1.63).  The $^{238}$U-late chain
  contributes $\sim$9$\times$ more than $^{232}$Th-late from detector
  components, because the $^{214}$Bi line at 2448~keV sits directly
  under $Q_{\beta\beta}$ while the $^{208}$Tl line at 2615~keV is
  partially rejected by single-scatter and veto cuts.

\item \textbf{Cavern walls} (11.6 counts, 33\%): the rock gamma flux
  is the single largest individual source.  Additional steel shielding
  (8~cm above the water tank) was added specifically for the
  $0\nu\beta\beta$ program.

\item \textbf{Internal backgrounds} (0.77 counts, 2\%): $^{222}$Rn
  (0.45), $^{137}$Xe from neutron capture (0.28), $2\nu\beta\beta$
  (0.01), $^{8}$B solar neutrinos (0.03).
\end{enumerate}

% ==================================================================
\section{Sensitivity}
% ==================================================================

Using a profile likelihood ratio analysis over the extended fiducial
volume (5.6~tonnes, 1360~kg$\cdot$years of $^{136}$Xe exposure) and
the energy range 2000--2700~keV:

\begin{itemize}
\item \textbf{Natural abundance} (8.9\% $^{136}$Xe): median 90\% CL
  exclusion sensitivity $T_{1/2}^{0\nu} > 1.06 \times 10^{26}$~years
  after 1000~live-days.

\item \textbf{90\% enrichment}: $T_{1/2}^{0\nu} > 1.06 \times
  10^{27}$~years after 1000~live-days (same detector, no other
  upgrades).

\item \textbf{Effective neutrino mass}: $m_{\beta\beta} = 53$--164~meV
  (range from nuclear matrix element uncertainty).
\end{itemize}

\subsection{Sensitivity dependence on key parameters}

\begin{itemize}
\item \textbf{Energy resolution}: modest impact from 1.0\% to
  $\sim$1.4\%; significant degradation above 2.0\% (the $^{208}$Tl
  peak leaks into the ROI).

\item \textbf{Vertex separation $\Delta z$}: strong dependence.
  Single-scatter rejection is the most powerful background cut for
  external gammas at these energies.  Below $\sim$2~mm, signal
  efficiency begins to drop (bremsstrahlung from the two electrons
  creates resolvable multi-site topologies).

\item \textbf{PLR vs.\ cut-and-count}: the profile likelihood
  analysis with extended volume and energy range improves sensitivity
  by a factor of $\sim$2 over a simple cut-and-count in the inner
  967~kg.
\end{itemize}

% ==================================================================
\section{LXeMC Source Model}
% ==================================================================

LXeMC models each background source as an independent geometric
object with associated material, mass, and radioactive activity.
Source geometry is defined in \texttt{source\_geometry\_lz\_v1.json},
separate from the tracking geometry used for photon transport.  The
two files describe the same physical detector but serve different
stages of the simulation: source-side emission and self-shielding
vs.\ per-step tracking inside the active LXe.

\subsection{Modeling categories}

Every source object belongs to one of two transport categories:

\begin{description}[leftmargin=2.0cm,style=nextline]
\item[\texttt{exact}] Real material with Klein--Nishina source-side
  transport.  The drawn geometry reproduces the measured mass within
  $\sim$1\%.  Used for Ti cryostat shells, PTFE reflector, and Ti
  field-cage rings.

\item[\texttt{mass\_equivalent}] Virtual carrier filled with vacuum
  (no source-side self-shielding).  The carrier geometry defines where
  gammas are emitted; the absolute activity is anchored by an explicit
  \texttt{equivalent\_mass\_kg} field rather than by geometric volume.
  Used for PMTs, PMT bases, PMT structures, cables, MLI insulation,
  FC resistors, TPC sensors, and grid holders.
\end{description}

\subsection{Activity convention}

The JSON fields \texttt{Bi214\_mBq\_per\_kg} and
\texttt{Tl208\_mBq\_per\_kg} store \emph{chain activities}, not
isotope-specific gamma rates:

\begin{itemize}
\item \texttt{Bi214\_mBq\_per\_kg} $=$ $^{238}$U-late chain activity
  ($^{226}$Ra onward), matching the paper's ``$^{238}$U-late'' column.

\item \texttt{Tl208\_mBq\_per\_kg} $=$ $^{232}$Th-late chain activity
  ($^{224}$Ra onward), matching the paper's ``$^{232}$Th-late'' column.
\end{itemize}

To obtain the relevant gamma rates for the $0\nu\beta\beta$ ROI:
\begin{align}
  \text{$^{214}$Bi 2448~keV rate} &= \text{Bi214\_mBq\_per\_kg}
    \times 0.015 \quad (\text{BR}) \\
  \text{$^{208}$Tl 2615~keV rate} &= \text{Tl208\_mBq\_per\_kg}
    \times 0.359 \times 0.998 \quad
    (\text{$^{212}$Bi$\to$$^{208}$Tl BR} \times \text{$\gamma$ BR})
\end{align}

\subsection{Cryostat sources}

Table~\ref{tab:cryostat} lists the six cryostat shells modeled in the
source geometry.  All shells are modeled as \texttt{exact} Ti
sources with Klein--Nishina self-shielding.  The geometric masses
reproduce the paper's component masses within 1\%; the sum of the six
drawn shells (1425~kg) is 55\% of the paper's total cryostat mass
(2590~kg), with the difference due to flanges, ports, and structural
elements not drawn as separate source objects.

\begin{table}[h]
\centering
\caption{Cryostat source shells.  All Ti, \texttt{exact}
  approximation, KN transport.  Activities: $^{238}$U-late $= 0.08$,
  $^{232}$Th-late $= 0.22$~mBq/kg (upper limits, uniform across all
  shells).}
\label{tab:cryostat}
\begin{tabular}{@{}llrrrl@{}}
\toprule
{Source} & {Shape} & {$R_\text{in}$ (cm)} & {$t$ (cm)} &
  {$h_\text{half}$ or AR} & {Geom.\ mass (kg)} \\
\midrule
OCV barrel  & cyl.\ shell & 90.8  & 0.70 & $h = 106.25$ & 384.0 \\
OCV top     & disk (2:1)  & 91.5  & 0.90 & AR $= 2.0$   & 147.3 \\
OCV bottom  & disk (2:1)  & 91.5  & 1.50 & AR $= 2.0$   & 245.4 \\
ICV barrel  & cyl.\ shell & 82.1  & 0.90 & $h = 94.915$ & 399.4 \\
ICV top     & disk (2:1)  & 83.0  & 0.80 & AR $= 2.0$   & 107.7 \\
ICV bottom  & disk (3:1)  & 83.0  & 1.20 & AR $= 3.0$   & 141.4 \\
\midrule
\multicolumn{5}{@{}l}{\textbf{Total (6 shells)}} & \textbf{1425} \\
\bottomrule
\end{tabular}
\end{table}

\subsection{Field-cage and internal sources}

Table~\ref{tab:fc_sources} lists the field-cage and internal detector
sources.  FC\_PTFE and FC\_rings are \texttt{exact} sources whose
drawn geometry reproduces the bb0nu reference masses (184~kg and
93~kg).  All other entries are \texttt{mass\_equivalent} virtual
carriers.

\begin{table}[h]
\centering
\caption{Field-cage and internal sources.  Dimensions refer to the
  full FC envelope ($z \in [-13.75, 145.6]$~cm, $h_\text{half} =
  79.675$~cm).}
\label{tab:fc_sources}
\begin{tabular}{@{}llrrrrrr@{}}
\toprule
{Source} & {Material} & {$R_\text{in}$} & {$t$} &
  {Mass} & {$^{238}$U} & {$^{232}$Th} & {Transport} \\
 & & {(cm)} & {(cm)} & {(kg)} & {(mBq/kg)} & {(mBq/kg)} & \\
\midrule
FC\_PTFE     & PTFE   & 72.8  & 1.147 & 184$^a$  & 0.04 & 0.01 & KN \\
FC\_rings    & Ti     & 74.3  & 0.277 & 93.0$^a$ & 0.35 & 0.24 & KN \\
FC\_resistors& Vacuum & 72.7  & 0.1   & 0.06$^b$ & 1350 & 2010 & transp.\ \\
FC\_sensors  & Vacuum & 72.6  & 0.1   & 5.02$^b$ & 5.82 & 1.88 & transp.\ \\
FC\_topgrid  & SS     & 72.8  & 7.5   & 44.6$^b$ & 2.63 & 1.46 & KN \\
FC\_botgrid  & SS     & 72.8  & 7.5   & 22.3$^b$ & 2.63 & 1.46 & KN \\
\bottomrule
\multicolumn{8}{@{}l}{\footnotesize $^a$ Geometric mass from drawn
  dimensions.  $^b$ \texttt{equivalent\_mass\_kg} (virtual carrier).} \\
\multicolumn{8}{@{}l}{\footnotesize Grid holders: JSON total
  66.8~kg vs.\ paper 89.1~kg; difference is grid wire mass not
  modeled separately.}
\end{tabular}
\end{table}

\subsection{PMT and cable sources}

Table~\ref{tab:pmt_sources} lists the PMT-related sources.  All are
\texttt{mass\_equivalent} virtual carriers with transparent transport.
Masses and activities match paper Table~I exactly.

\begin{table}[h]
\centering
\caption{PMT and cable sources.  All virtual carriers (Vacuum,
  transparent transport).  ``top+bot'' entries are modeled as separate
  disk-shaped sources at their respective $z$ positions; the table
  shows per-source values.}
\label{tab:pmt_sources}
\begin{tabular}{@{}llrrr@{}}
\toprule
{Source} & {Shape} & {Mass (kg)} &
  {$^{238}$U (mBq/kg)} & {$^{232}$Th (mBq/kg)} \\
\midrule
PMT\_TOP\_PMTs       & cylinder & 47.07 & 3.22 & 1.61 \\
PMT\_BOT\_PMTs       & cylinder & 44.83 & 3.22 & 1.61 \\
\quad \emph{sum}     &          & \emph{91.90} & & \\[2pt]
PMT\_TOP\_bases      & cylinder &  1.434 & 75.9 & 33.1 \\
PMT\_BOT\_bases      & cylinder &  1.366 & 75.9 & 33.1 \\
\quad \emph{sum}     &          & \emph{2.80} & & \\[2pt]
PMT\_TOP\_structure  & cylinder & 83.0  & 1.60 & 1.06 \\
PMT\_BOT\_structure  & cylinder & 83.0  & 1.60 & 1.06 \\
\quad \emph{sum}     &          & \emph{166.0} & & \\[2pt]
PMT\_BARREL\_cables  & cyl.\ shell & 88.7 & 4.31 & 0.82 \\
\midrule
\multicolumn{5}{@{}l}{\emph{Skin PMTs (all share 46.0 / 14.9 mBq/kg):}} \\
PMT\_BARREL\_R8520       & cyl.\ shell & 6.10  & 46.0 & 14.9 \\
PMT\_BARREL\_R8778\_lower & cyl.\ shell & 1.31  & 46.0 & 14.9 \\
PMT\_BOT\_R8778\_dome    & cylinder    & 1.18  & 46.0 & 14.9 \\
\quad \emph{sum}         &             & \emph{8.59} & & \\
\bottomrule
\end{tabular}
\end{table}

\subsection{Insulation source}

The MLI cryostat insulation is modeled as a virtual cylindrical shell
on the ICV outer wall ($R_\text{in} = 83.0$~cm, $t = 1.0$~cm,
$h_\text{half} = 94.92$~cm).  Material is Vacuum (transparent
transport); \texttt{equivalent\_mass\_kg} $= 13.8$~kg.  Activities:
$^{238}$U-late $= 11.1$, $^{232}$Th-late $= 7.79$~mBq/kg.

\subsection{Components not modeled}

Two paper Table~I entries are not present in the source geometry JSON:

\begin{itemize}
\item \textbf{Outer detector system} (22\,900~kg, 4.71 / 3.73~mBq/kg,
  2.79~cts/1000d): the Gd-doped liquid scintillator and acrylic
  vessels.  These are external to the LXe and currently not modeled as
  gamma sources.

\item \textbf{Other components} (438~kg, 1.83 / 1.65~mBq/kg,
  2.41~cts/1000d): miscellaneous small items.
\end{itemize}

Together these contribute 5.2~cts/1000d (15\% of the total
background).  Their omission is conservative for source-flux
validation but should be noted when comparing absolute background
rates to the paper.

\subsection{Tracking geometry}

The tracking geometry (\texttt{detector\_lz\_v3.json}) models the
detector interior as seen by the per-step photon transport.
Table~\ref{tab:tracking} summarizes the key dimensions.

\begin{table}[h]
\centering
\caption{Tracking geometry dimensions.  All regions are coaxial with
  the detector axis.  $z = 0$ is the cathode; $z = 145.6$~cm is the
  gate.}
\label{tab:tracking}
\begin{tabular}{@{}llrrl@{}}
\toprule
{Region} & {Shape} & {$R$ or $[R_\text{in}, R_\text{out}]$} &
  {$z$ range} & {Tag} \\
 & & {(cm)} & {(cm)} & \\
\midrule
LZ\_detector & domed container & 82.1 &
  $[-41.3, 195.5]^*$ & vacuum \\
AirDome      & cap (up)        & 82.1 &
  $[145.6, 186.7]$    & vacuum \\
LXe\_passive & capped cyl.     & 82.1 &
  $[-41.3, 145.6]$    & passive\_lxe \\
\midrule
TopActive    & cylinder & 72.8 &
  $[96.0, 145.6]$ & tpc\_active \\
BarrelActive & cyl.\ shell & $[39.0, 72.8]$ &
  $[26.0, 96.0]$  & tpc\_active \\
BottomActive & cylinder & 72.8 &
  $[0.0, 26.0]$   & tpc\_active \\
FV           & cylinder & 39.0 &
  $[26.0, 96.0]$  & fv \\
\midrule
FC\_PTFE     & cyl.\ shell & $[72.8, 73.95]$ &
  $[-13.75, 145.6]$ & structural \\
FC\_rings    & cyl.\ shell & $[74.3, 74.58]$ &
  $[-13.75, 145.6]$ & structural \\
Skin         & cyl.\ shell & $[74.58, 82.1]$ &
  $[-13.75, 145.6]$ & skin \\
\bottomrule
\multicolumn{5}{@{}l}{\footnotesize $^*$ Includes top and bottom
  ellipsoidal dome caps (AR $= 2$ and $3$).}
\end{tabular}
\end{table}

Key consistency checks against the paper:

\begin{itemize}
\item \textbf{TPC drift}: $z \in [0, 145.6]$~cm, fully tiled by
  BottomActive + FV + BarrelActive + TopActive.  Matches paper's
  145.6~cm drift length.

\item \textbf{TPC radius}: 72.8~cm (= 145.6/2~cm inner diameter).
  Matches paper.

\item \textbf{FV}: $r < 39$~cm, $z \in [26, 96]$~cm, $\sim$967~kg.
  Matches paper Section~III.

\item \textbf{Reverse field region}: FC structures extend to $z =
  -13.75$~cm, consistent with the paper's 13.75~cm RFR depth below
  the cathode.

\item \textbf{ICV inner radius}: LZ\_detector envelope at $R =
  82.1$~cm, matching ICV\_barrel $R_\text{in}$ in the source JSON.

\item \textbf{Skin}: modeled as a uniform 7.5~cm shell.  The paper
  describes a position-dependent thickness (4~cm at top, 8~cm at
  cathode level); the JSON uses a simplified constant thickness that
  provides the correct radial partition.

\item \textbf{FC dimensions}: identical in source and tracking JSON
  (validated by the test suite).
\end{itemize}

% ==================================================================
\section{Key Numbers for LXeMC Simulation}
% ==================================================================

Table~\ref{tab:lxemc} collects the parameters adopted by LXeMC from
the LZ $0\nu\beta\beta$ analysis.

\begin{table}[h]
\centering
\caption{Parameters adopted by LXeMC from the LZ $0\nu\beta\beta$
  analysis.}
\label{tab:lxemc}
\begin{tabular}{@{}lll@{}}
\toprule
Parameter & Value & Use in LXeMC \\
\midrule
FV dimensions & $r < 39$~cm, $26 < z < 96$~cm & \texttt{FVGeometry} \\
Energy resolution & $\sigma/E = 1.0\%$ at $Q_{\beta\beta}$ & event smearing \\
$\Delta z$ clustering & 3~mm & SS/MS classification \\
Skin veto threshold & 100~keV & \texttt{veto\_skin} \\
TPC veto threshold & 10~keV & \texttt{veto\_TPC} \\
ROI & $\pm 1\sigma$ around $Q_{\beta\beta}$ & event selection \\
Component masses & Tables~\ref{tab:cryostat}--\ref{tab:pmt_sources} & source geometry JSON \\
Component activities & Tables~\ref{tab:cryostat}--\ref{tab:pmt_sources} & source flux generation \\
Tracking dimensions & Table~\ref{tab:tracking} & tracking geometry JSON \\
\bottomrule
\end{tabular}
\end{table}

\end{document}
